\section{Cutting to obtain separation}\label{sec:cutting}

The surface estimate controls the size of a closed stable weighted minimal
surface, but does not make an arbitrary separating construction connected.
We therefore first remove the topological obstruction: after cutting along
a locally finite family of small stable surfaces, every compact connected
two-sided hypersurface neatly embedded in a resulting piece will separate
it. Throughout this
section, $M$ is complete and $D(g,t)\geq1$.

A hypersurface in a manifold with boundary is \emph{neatly embedded} if
its boundary is exactly its intersection with the ambient boundary and
the intersection is transverse. A closed surface means a compact surface
without boundary. All hypersurfaces considered below are smooth.

\begin{proposition}\label{prop:separating-cut}
There is a locally finite family $\mathcal S$ of pairwise disjoint closed
connected embedded two-sided stable weighted minimal surfaces in $M$
such that the following statements hold. Let $N$ be the manifold obtained
by cutting $M$ along $\mathcal S$, and let $q:N\to M$ be the natural
gluing map.
\begin{enumerate}
\item The map $q$ is proper and is a quotient map. Each connected component
$P$ of $N$, with the intrinsic length metric of the pulled-back metric
$g$, is proper.
\item Each boundary component of $P$ is a sphere or a real projective
plane, is stable weighted minimal, and has intrinsic $g$-diameter at most
$2\pi$. Each surface of $\mathcal S$ gives two boundary copies in $N$.
\item Every compact connected neatly embedded two-sided hypersurface
in $P$ separates $P$.
\end{enumerate}
\end{proposition}

\subsection{A maximal locally finite family}

Order by inclusion the families of pairwise disjoint closed connected
embedded two-sided stable weighted minimal surfaces for which no finite
subfamily disconnects $M$. The empty family is admissible, and the union
of a chain is admissible because each finite subfamily lies in one member
of the chain. Zorn's lemma gives a maximal family $\mathcal S$.

We first verify local finiteness. By Proposition~\ref{prop:stable-surface},
every member has intrinsic, and hence ambient, diameter at most $2\pi$.
Consequently, all members meeting a fixed compact set $A\subset M$ lie
in the compact closed $2\pi$-neighborhood of $A$. Here completeness is
used through the properness of the Riemannian distance on $M$.

For any finite subfamily $S_1,\ldots,S_k$, the classes
$[S_i]\in H_2(M;\Z/2)$ are linearly independent. Indeed, the connected
complement of the subfamily contains a path joining the two sides of
$S_i$; closing it by a short transverse arc gives a loop meeting $S_i$
once and missing the other surfaces. Modulo two, intersection with this
loop pairs to one with $[S_i]$ and to zero with the other classes. This
uses only the intersection pairing of compact cycles, so neither
orientability nor compactness of $M$ is required. A compact smooth
submanifold neighborhood of the closed $2\pi$-neighborhood of $A$ has
finite-dimensional $H_2(-;\Z/2)$. Its image in $H_2(M;\Z/2)$ therefore
bounds the size of every finite subfamily meeting $A$. Thus only finitely
many members meet $A$.

Local finiteness makes cutting well-defined. In a tubular chart at a cut
surface, $q$ glues the two closed half-charts along their boundary; away
from the cuts it is a diffeomorphism. These charts show that $q$ is a
quotient map and is proper locally over relatively compact neighborhoods.
A finite such cover of a compact subset of $M$ proves global properness.
The pulled-back metric induces the manifold topology, also for its
intrinsic length metric on each component $P$. A closed bounded subset
of $P$ maps into a bounded subset of $M$; it is closed in the compact
inverse image of a sufficiently large closed ball. It is therefore
compact. This proves properness of $P$.

The boundary assertions in Proposition~\ref{prop:separating-cut} now
follow from Proposition~\ref{prop:stable-surface}. It remains to prove
separation. Fix a connected component $P$ of $N$. We first treat closed
hypersurfaces in the interior; this is
the only step requiring minimization.

\subsection{Producing an additional stable surface}

\begin{lemma}\label{lem:closed-separation}
Every closed connected embedded two-sided hypersurface in
$\operatorname{int}P$ separates $P$.
\end{lemma}

\begin{proof}
Suppose that $F\subset\operatorname{int}P$ is such a hypersurface and
does not separate. Fix a coorientation of $F$. A loop crossing $F$ once
has nonzero algebraic crossing number. If it reverses ambient orientation,
traverse it twice. A general-position perturbation in dimension three
then gives a smooth embedded, transverse, orientation-preserving loop
$\gamma\subset\operatorname{int}P$ with nonzero algebraic crossing
number with $F$. The crossing number is unchanged by this perturbation.

\paragraph{A compact minimization problem.}
Let $\pi:\widehat P\to P$ be the orientation double cover, including
both copies when $P$ is orientable. Give $\widehat P$ its canonical
orientation and write $\tau$ for the deck involution. The coorientation
of $F$ and the ambient orientation orient its full lift as an integral
$2$-cycle $T_F$ satisfying
\[
\tau_\#T_F=-T_F.
\]
Since $\gamma$ preserves orientation, it has two closed lifts. In a small
tubular neighborhood of one of them, choose a smooth compactly supported
closed $2$-form $\beta$ such that $T_F(\beta)\ne0$. Concretely, the
normal disk bundle of this lifted loop is an oriented trivial $2$-disk
bundle. A disk area form with a smooth bump of integral one, extended
along the loop and by zero outside the tube, is a Thom form. Its
evaluation is the signed crossing number, up to the orientation chosen
for the disk. The tube and $\supp\beta$ lie in the interior.
Preserving this nonzero evaluation will detect a nonseparating component,
once the surface estimate has made the components meeting the form
support compact.

Choose a relatively open set $W\subset P$ with compact closure that
contains $F$ and the closed $10$-neighborhood of
$\pi(\supp\beta)$ in the intrinsic metric of $P$. Such a set exists
by properness. Let $J\subset P$ be a compact set containing $\overline W$.
In the conformal metric
\begin{equation}\label{eq:area-metric}
\bar g=e^{-t}g,
\end{equation}
surface area is precisely area weighted by $e^{-t}$. Minimize
$\bar g$-mass among integral $2$-cycles $T$ supported in $\pi^{-1}(J)$
that satisfy
\begin{equation}\label{eq:current-constraints}
\tau_\#T=-T,
\qquad T(\beta)=T_F(\beta).
\end{equation}
These are absolute cycles, not relative cycles. To interpret them at
the boundary, smoothly extend $P$ and $\bar g$ across $\partial P$,
take the orientation cover of that extension, and regard the currents
as currents in the extension with the indicated support constraint.
The form $\beta$ extends by zero. The competitor $T_F$ shows that the
class is nonempty. Federer--Fleming compactness and closure
\cite{FedererFleming1960}, followed
by lower semicontinuity of mass, give a minimizer $T_0$: support is
confined to a fixed compact set, the boundaries vanish, and both
conditions in \eqref{eq:current-constraints} are closed under weak
convergence.

In a sufficiently small ball over $W\cap\operatorname{int}P$, $T_0$
is locally mass minimizing. A compactly supported integral-cycle change
there evaluates to zero on $\beta$, since $\beta$ is exact in the ball.
Making the opposite change in the disjoint deck translate preserves
\eqref{eq:current-constraints}; a strict decrease would give a strict
decrease of twice the amount upstairs. We must establish the same
minimizing property across the boundary before using interior regularity.

\paragraph{Comparison across a minimal boundary.}
Each boundary component is minimal for $\bar g$, because its
$\bar g$-area is its original weighted area. Near any point of such a
component there is a foliation by $\bar g$-minimal disks, with the
boundary disk as one leaf. Here is the local construction, which does
not require a globally stable extension of the boundary. On a sufficiently
small disk in that leaf, the Dirichlet Jacobi operator, with the stability
sign, has strictly positive first eigenvalue: its zeroth-order term is
bounded and the Poincar\'e inequality dominates it after shrinking the
disk. Solve the minimal-graph equation with small constant normal boundary
height by the implicit function theorem, using the standard Dirichlet
elliptic theory \cite{GilbargTrudinger2001}. The derivative at height zero
solves the Jacobi equation with boundary value one and is strictly
positive by the maximum principle. After shrinking the disk and the
height interval, these graphs give the asserted foliation.

Let $Y$ be its unit normal field, directed toward the allowed side of
the boundary. Since all leaves are minimal,
$\Div_{\bar g}Y=0$. In an oriented flow box the $2$-form
\[
\alpha=\iota_Y\,d\vol_{\bar g}
\]
is closed, annihilates $Y$, and is invariant under its flow. Projection
along this flow onto the boundary leaf therefore pulls its area form
back to $\alpha$. As $\alpha$ has comass one, this projection has
$2$-Jacobian at most one on every tangent $2$-plane.

To apply this to currents, choose a contractible flow box $U'$ upstairs,
with coordinates $(z,s)$, whose boundary leaf is $s=0$ and whose allowed
side is $s\geq0$. Require that the allowed side lie over $W$, and that
$\overline{U'}$ and $\tau\overline{U'}$ be disjoint. Take a smaller
box $U\Subset U'$. Suppose a compactly supported integral $2$-cycle
$S$ in $U$ gave a strictly smaller mass for $T'=T_0+S$. Use the
locally Lipschitz map
\[
R(z,s)=(z,\max\{0,s\})
\]
near $\overline U$, interpolating to the identity before leaving $U'$,
and keeping $R$ equal to the identity everywhere on the allowed side.
The interpolation can be made in the flow coordinate and can be taken
to map $U'$ into itself. Every portion of $T'$ on the exterior side is
supported in $U$, where the exact projection rule holds. On the allowed
side, including tangent planes carried by the boundary leaf, $R$ is
the identity. The area formula and the preceding Jacobian bound give
\begin{equation}\label{eq:boundary-retraction-mass}
\mathbf M_{\bar g}(R_\#T')\leq \mathbf M_{\bar g}(T').
\end{equation}
No Jacobian bound for the interpolation region is needed, since $T'$
has no exterior mass there.

Put $S'=R_\#T'-T_0$. Pushforward preserves the cycle condition, and
$R_\#T_0=T_0$. Thus $S'$ is an integral cycle supported on the allowed
side in a compact subset of $U'$. Both $S'$ and $\tau_\#S'$ evaluate
to zero on $\beta$: use exactness in the respective contractible boxes,
with primitives cut off beyond the cycle supports. Consequently
\[
T_0+S'-\tau_\#S'
\]
is again supported in $\pi^{-1}(J)$ and satisfies
\eqref{eq:current-constraints}. The boxes are disjoint, and antisymmetry
of $T_0$ makes the two mass changes equal. By
\eqref{eq:boundary-retraction-mass} their sum is strictly negative,
contradicting minimality. This proves local mass minimization in full
neighborhoods across the boundary.

\paragraph{Regular sheets and their coorientations.}
We now use codimension-one interior regularity for mass-minimizing
integral currents: in a smooth Riemannian $3$-manifold, a locally
mass-minimizing integral $2$-current without boundary is, near each
support point, a smooth embedded minimal surface with locally constant
integer multiplicity and a smooth orientation
\cite[Chapter~7, Section~5]{Simon2014}. The full-neighborhood
comparison just proved supplies precisely the interior hypothesis even
at points of $\partial P$. The strong maximum principle implies that
any sheet touching the boundary agrees there with the boundary leaf.
Hence, over $W$, the projected support is smooth and each connected
component is either contained in $\partial P$ or disjoint from it.

In a sufficiently small evenly covered neighborhood of a projected
support point, regularity gives an embedded sheet in each lift. Deck
invariance identifies their projections, so the projected support is
itself an embedded surface. Moreover, the orientations
give coorientations downstairs: $\tau$ reverses both the current
orientation and the ambient orientation, so preserves the resulting
normal direction. The absolute multiplicities on the two lifts of a
connected projected component agree. Every component in
$W\cap\operatorname{int}P$ is therefore a two-sided stable weighted
minimal surface, initially possibly noncompact. To see the asserted
stability, extend a compactly supported normal variation of one sheet
to an ambient flow isolating that sheet and supported in
$W\cap\operatorname{int}P$. Its lifts preserve antisymmetry. They
preserve evaluation on $\beta$ by the homotopy formula, because
$d\beta=0$ and $\partial T_0=0$. Minimality of $T_0$, and the locally
constant multiplicities, then give the weighted stability inequality
for that sheet.

We have obtained stable surfaces near the detecting form. The next
point is that the surface length estimate makes the relevant components
compact, even though no completeness of these intermediate sheets has
been asserted.

\paragraph{The components detected by the form are closed.}
Let $\Sigma$ be a component of the projected support on $W$ meeting
$\pi(\supp\beta)$, and choose $x$ in that intersection. This component
is disjoint from $\partial P$. Its intersection with
$\overline B_P(x,10)$ is compact. Indeed, this ball is compact and is
contained in $W$; the projected support is closed relative to $W$
because the cover is finite, and its components are closed there.
The regular embedded surface structure is the subspace topology.

If $\Sigma$ were noncompact, choose a connected relatively compact
smooth domain $\Omega_0\subset\Sigma$ containing this intersection
in its interior. Its nonempty boundary lies outside
$\overline B_P(x,10)$. Choose another connected relatively compact
smooth surface domain $\Omega_1$ with
$\overline{\Omega_0}\subset\Omega_1$. Stability holds throughout
$\Omega_1$, and $D(g,t)\geq1$. Lemma~\ref{lem:stable-exit}, with
$\mu=0$ and $c=1$, produces a path from $x$ to
$\partial\Omega_0$ of induced $g$-length at most $2\pi$. Its image
in $P$ has no greater length. This contradicts the distance of
$\partial\Omega_0$ from $x$, which is greater than $10$. Thus
$\Sigma$ is a closed surface in the interior of $P$.

Only finitely many projected components meet the form support, by
compactness and local regularity. Since $T_0(\beta)\ne0$, at least
one of these closed components is nonseparating in $P$. For otherwise,
for each component choose the closure of one of its sides, with the
choice agreeing with its coorientation. On the orientation cover,
Stokes' theorem gives zero for the integral of $\beta$ over its full
oriented lift. This remains valid when the side is noncompact, since
$\beta$ is compactly supported; boundary portions in $\partial P$
contribute nothing, since its support lies in the interior. Multiplying
by the constant multiplicities and summing would give $T_0(\beta)=0$.

Let $S$ be such a nonseparating component. It is closed, connected,
embedded, two-sided, stable weighted minimal, and lies in
$\operatorname{int}P$. Its image in $M$ is disjoint from every member
of $\mathcal S$. Adding this image preserves the defining property of
$\mathcal S$. In fact, after deleting any finite subfamily the ambient
manifold remains connected, and the two sides of $S$ can still be joined
while avoiding $S$: join them in $P\setminus S$ and push the path into
the interior of $P$. A connected two-sided hypersurface has at most two
complementary components, so this joining path means that the new
surface does not disconnect that complement. We have contradicted
maximality, proving the lemma.
\end{proof}

\subsection{From closed to neat hypersurfaces}

We finish the proof of Proposition~\ref{prop:separating-cut}. A compact
connected two-sided neat hypersurface has a bicollar, and its complement
has at most two components. If it does not separate, the bicollar and
a path joining its two sides produce a loop in the interior crossing it
once. Suppose $F$ is such a nonseparating neat hypersurface. Its bicollar
defines a smooth map
\[
u:P\longrightarrow S^1
\]
which is constant outside a compact set, goes once around the circle
along the transverse collar coordinate, and has $F$ as a regular level.
The map has nonzero degree on a loop crossing $F$ once. One may use a
transverse profile that is constant near both ends of the bicollar, so
extension by the same constant value outside the collar is smooth.

Every boundary component of $P$ has finite fundamental group. The
restriction of $u$ to a collar of any such component therefore has a
smooth real lift under $\R\to S^1$. Only finitely many boundary
components meet the compact region where $u$ is not constant, by local
finiteness. On their collars, multiply the difference between this lift
and a chosen lift of the outside constant by a smooth cutoff that is
zero near $\partial P$ and one near the interior end. Exponentiating
gives a compactly supported homotopy of $u$ to a smooth map $u_1$ that
equals the same constant on a neighborhood of all of $\partial P$.
Its degree on the crossing loop is unchanged.

Choose a regular value of $u_1$ distinct from the constant value.
Its inverse image is a compact cooriented hypersurface in
$\operatorname{int}P$, with finitely many connected components.
Perturb the crossing loop to be transverse to it. The sum of its signed
crossings with these components is the nonzero degree of $u_1$ on the
loop. If every component separated $P$, every such signed crossing
number would be zero. Thus a component would contradict
Lemma~\ref{lem:closed-separation}. This proves the final assertion of
Proposition~\ref{prop:separating-cut}.
